%% notes-tex/genesis/genesis-proposal-redacted/sections/background.tex
%% [[genesis-proposal-redacted]]
%% https://github.com/Mjvolk3/torchcell/tree/main/notes-tex/genesis/genesis-proposal-redacted/sections/background.tex
%% SOURCE: hand-transcribed from the proposal docx, "Background/Introduction".
%% Group names replaced by institutional roles. Figure 1 is the image embedded in
%% the docx (word/media/image1.png), converted to PDF unchanged.

\section{Background and introduction\secstatus{todo}}\label{sec:background}

Metabolic engineering has been increasingly used to produce chemicals, fuels, and
materials\textsuperscript{1}. Developing AI foundation models for metabolic
engineering is critical for addressing the complexity and scale of biological design
problems. Engineering microbial metabolism requires navigating an enormous
combinatorial space of genes, enzymes, regulatory interactions, and environmental
conditions. Traditional modeling and machine-learning approaches are typically
trained for narrow tasks and rely on limited labeled datasets, making them difficult
to generalize across organisms and pathways. In contrast, foundation models -- large
AI models pretrained on diverse biological datasets such as genomic sequences,
metabolic networks, protein structures, and multi-omics measurements -- can learn
transferable representations of biological systems and be adapted to many downstream
tasks with relatively little additional data\textsuperscript{2}. In essence, by
integrating heterogeneous biological data and capturing complex dependencies within
cellular systems, AI foundation models link genotype and phenotype (Focus area 2B)
and can guide pathway discovery, retrosynthetic planning, and strain optimization
while reducing costly experimental trial-and-error\textsuperscript{3}. Ultimately,
they could transform metabolic engineering into a more predictive and scalable
discipline, enabling rapid development of sustainable biomanufacturing platforms for
renewable chemicals, fuels, and materials (AI advantages).

The three partner groups (the lead university group, the LBNL group, and the ANL
group) are pioneers in developing AI models for metabolic engineering. For example,
the lead group developed an AI-driven automated workflow for pathway engineering in
2019\textsuperscript{4} while the LBNL group developed the Automated Recommendation
Tool (ART), a tool that leverages machine learning and probabilistic modeling
techniques to guide metabolic engineering in a systematic fashion in
2020\textsuperscript{5}. Moreover, the ANL group developed GenSLM, a DNA large
language model (LLM), which was able to predict novel variants in the newly emerging
coronavirus genomes (SARS-CoV-2) in 2023. Recently, in an unpublished study, the lead
group developed TorchCell, a framework to facilitate building metabolic engineering
foundation model (MEFM) with specific applications in \org{Saccharomyces cerevisiae},
a preeminent industrial platform organism (Fig.~\ref{fig:one}). Yeast genes are
embedded using various pretrained DNA and protein embedding models, and
protein-protein interactions, regulatory networks, and metabolism are then used to
build adjacency matrices among genes, reactions, and metabolites demonstrating
multimodal input representation of metabolic engineering experiments. We select
gene-gene epistasis, commonly referred to as gene-gene interactions (GGI), as our
supervised prediction label since it is a well-defined representation of the
nonlinear complexity of biological systems that makes it difficult to reliably
engineer commercially viable strains. The TorchCell Cell Graph Transformer (CGT)
model is trained on a dataset queried from the TorchCell Database, which aggregates
multiple original datasets (Fig.~\ref{fig:one}a--c). It achieves a Pearson
correlation of $0.454 \pm 0.006$ and a Spearman correlation of $0.421 \pm 0.004$,
with increased noise at the extremes of the GGI range, as expected
(Fig.~\ref{fig:one}d). We demonstrate state-of-the-art performance from benchmarking
against flux balance analysis (FBA), the common mechanistic, constraint-based
modeling method, and two recent deep learning frameworks DCell and
DANGO\textsuperscript{9--11} (Fig.~\ref{fig:one}e). Notably, many foundation models
have been developed for protein engineering; however, no foundation model has been
developed for metabolic engineering. This is largely due to the fragmentation of
phenotypic data across diverse experimental datasets, where individual studies are
typically too small to support deep learning approaches, preventing the field from
understanding how models scale across metabolic engineering data types. TorchCell
fills this gap by aggregating genotype, phenotype, and environmental data relevant
to metabolic engineering.

\tcfigfit[0.98]{figures/genesis-fig1.pdf}{TorchCell database construction and
\org{S. cerevisiae} application predicting cellular epistasis. \textbf{a},
Literature and public databases are annotated with TorchCell's experiment ontology,
enabling coherent querying across all data in the knowledge graph database.
Experiments are added to the TorchCell Graph database hosted by National Center for
Supercomputing Applications (NCSA), which can be queried to construct datasets large
enough to enable DL applications. \correction{Correction (2026-10-22): the
morphology and expression values shown in panel a as held at submission are now
known to be in error.} \textbf{b}, TorchCell's minimal requirements for
genomic representation are a full genome sequence and a genome annotation file, but
layers of additional graphical knowledge like regulatory networks and metabolism can
be added if annotations exist. Full representations of experiments are often highly
repetitive, so TorchCell compresses them into perturbation operators over a reference
state, facilitating training on millions of strains. \textbf{c}, Cell Graph
Transformer (CGT) takes multimodal cell inputs, embeds them into cell latent space
via a perturbation operator as described in recent virtual cell literature, and the
perturbed cell embeddings are decoded into different supervised labels, in this
case, GGI\textsuperscript{7--8}. \textbf{d}, TorchCell CGT validation on GGI.
\textbf{e}, TorchCell CGT achieves state-of-the-art performance in predicting
trigenic GGI in yeast compared to Yeast9 FBA, DCell, and Dango. \textbf{f},
TorchCell can be used to learn from previous strain experiments and help design the
next batch of strains for experimentation in autonomous labs like the NSF
iBioFoundry at UIUC.}{fig:one}
